\section{開發問題與解決方法}

\subsection{GLUT Callback 與多視窗狀態}

FreeGLUT callback 不會附帶 C++ object pointer，同一組 static callback 又可能由不同視窗觸發。程式因此維護 GLUT window ID 到 \lstinline|Window*| 的查找表，每次先辨識來源再轉交對應物件；Menu 也以 ID 對應實際動作。這層 routing 讓後續 Event、Element 與 Document 不必依賴全域的「目前視窗」，完整生命週期見附錄~\ref{app:gui}。

\subsection{Draft 與 Undo／Redo}

拖曳需要即時預覽，但若每次 Motion 都修改 Scene，Undo history 會累積大量中間狀態。程式將未完成圖形保留在 Tool 的 draft，由 Renderer 畫成 overlay；只有操作完成時，Document 才建立 Command 並修改 Scene。繪圖中按 Undo 會先取消 draft，沒有 draft 時才復原上一筆 Command，使預覽、正式資料與歷史紀錄各自有清楚的生命週期。

\subsection{粗線轉角的幾何例外}

自行產生粗線 triangle 時，尖角會讓外側交點離頂點過遠，接近 180 度折返或短 segment 也可能造成不穩定交點。程式先移除重複點，再以 cross product 與 dot product 區分情況；MITER 超過 miter limit 時退回 BEVEL，內側交點也設距離上限。求交失敗時則使用不依賴遠端交點的 BEVEL 或 ROUND 畫法，避免快速繪製短折線時出現尖刺與交錯 triangle。

\subsection{UTF-8 與 GFNT 資料}

FreeGLUT 的一般 keyboard callback 一次只提供一個 byte，若直接用 byte index 編輯，Backspace 可能切斷中文字。輸入框因此以 code point vector 保存文字，顯示時才重新編碼為 UTF-8。GFNT Loader 則在載入時檢查 Header、Glyph Table、Unicode 範圍、排序與 bitmap 邊界，render loop 只需 binary search glyph；格式與驗證流程見附錄~\ref{app:gfnt}。
